% GENERATED FILE. Edit canonical lessons, then run npm run generate:books.
% canonical-source-hash: fnv1a64:e7a8473b2d48076f
% canonical-lessons: GE-C461-nieselregen, GE-C461-nieseln, GE-C461-schauer, GE-C461-aufheiterung, GE-C461-aufklaren, GE-R461-review-454-461

\chapter{Schauer und Nieselregen}
\label{ch:ge-schauer-und-nieselregen}
\hlchaptermodality{461}

\begin{chapteropening}
\textbf{By the end of this chapter:} \emph{I can talk about light rain: drizzle, drizzling, the shower, the sunny spell and the sky clearing.}

Complete the last lesson of chapter 461: I can talk about light rain: drizzle, drizzling, the shower, the sunny spell and the sky clearing.
\end{chapteropening}

\section[der Nieselregen]{der Nieselregen — the drizzle}

\label{lesson:GE-C461-nieselregen}



\begin{quote}
Before the new one: say the German for overcast, then the German for to change suddenly.
\end{quote}

\subsection*{You'll want to know: der Nieselregen}
\textbf{der Nieselregen} — \textquotedblleft{}the drizzle\textquotedblright{}.

\begin{quote}
Bei Nieselregen gehen wir trotzdem spazieren.
\end{quote}

\begin{quote}
\emph{We go for a walk in the drizzle anyway.}
\end{quote}

No plural.

\begin{cousinweb}
\emph{Nieselregen} is built from \emph{nieseln} 'to drizzle' and \emph{Regen} 'rain', an inherited word and a cousin of English \emph{rain} (EtymDB).
\end{cousinweb}

\begin{grammarlens}[title={What you've built}]
One more word for telling, explaining and describing at B1.
\end{grammarlens}

\subsection*{Your turn}
\begin{itemize}
\raggedright
  \item \emph{Say it:} \emph{der Nieselregen}
  \item \emph{Say it:} \emph{der Nieselregen}, once more
  \item \emph{Say it:} a sentence of your own with \emph{der Nieselregen}
  \item \emph{Recall:} say the German for overcast, then the German for to change suddenly, then say \emph{der Nieselregen} again
\end{itemize}

\begin{tcolorbox}[breakable,colback=teal!4,colframe=teal!35!black,title={Before you move on}]
What does der Nieselregen mean? (\textquotedblleft{}The drizzle\textquotedblright{}.) Say it once more.
\end{tcolorbox}
\section[nieseln]{nieseln — to drizzle}

\label{lesson:GE-C461-nieseln}



\begin{quote}
Before the new one: say the German for to change suddenly, then the German for the drizzle.
\end{quote}

\subsection*{You'll want to know: nieseln}
\textbf{nieseln} — \textquotedblleft{}to drizzle\textquotedblright{}.

\begin{quote}
Es nieselt schon den ganzen Vormittag.
\end{quote}

\begin{quote}
\emph{It has been drizzling all morning.}
\end{quote}

Mostly with \emph{es}: \emph{es nieselt}, \emph{es hat genieselt}.

\begin{cousinweb}
The origin of \emph{nieseln} is not settled; EtymDB lists it without a source. It is sometimes linked to \emph{niesen} 'to sneeze', for the fine spray.
\end{cousinweb}

\begin{grammarlens}[title={What you've built}]
One more word for telling, explaining and describing at B1.
\end{grammarlens}

\subsection*{Your turn}
\begin{itemize}
\raggedright
  \item \emph{Say it:} \emph{nieseln}
  \item \emph{Say it:} \emph{nieseln}, once more
  \item \emph{Say it:} a sentence of your own with \emph{nieseln}
  \item \emph{Recall:} say the German for to change suddenly, then the German for the drizzle, then say \emph{nieseln} again
\end{itemize}

\begin{tcolorbox}[breakable,colback=teal!4,colframe=teal!35!black,title={Before you move on}]
What does nieseln mean? (\textquotedblleft{}To drizzle\textquotedblright{}.) Say it once more.
\end{tcolorbox}
\section[der Schauer]{der Schauer — the shower (of rain); a shiver}

\label{lesson:GE-C461-schauer}



\begin{quote}
Before the new one: say the German for the drizzle, then the German for to drizzle.
\end{quote}

\subsection*{You'll want to know: der Schauer}
\textbf{der Schauer} — \textquotedblleft{}the shower (of rain); a shiver\textquotedblright{}.

\begin{quote}
Am Nachmittag gibt es immer wieder kurze Schauer.
\end{quote}

\begin{quote}
\emph{In the afternoon there are short showers again and again.}
\end{quote}

Same form in the plural. \emph{Ein Schauer} is also a shiver.

\begin{cousinweb}
\emph{Schauer} is inherited through Middle High German \emph{schure} from Old High German \emph{scūr} 'storm, shower' (EtymDB); English \emph{shower} is a cousin.
\end{cousinweb}

\begin{grammarlens}[title={What you've built}]
One more word for telling, explaining and describing at B1.
\end{grammarlens}

\subsection*{Your turn}
\begin{itemize}
\raggedright
  \item \emph{Say it:} \emph{der Schauer}
  \item \emph{Say it:} \emph{der Schauer}, once more
  \item \emph{Say it:} a sentence of your own with \emph{der Schauer}
  \item \emph{Recall:} say the German for the drizzle, then the German for to drizzle, then say \emph{der Schauer} again
\end{itemize}

\begin{tcolorbox}[breakable,colback=teal!4,colframe=teal!35!black,title={Before you move on}]
What does der Schauer mean? (\textquotedblleft{}The shower (of rain); a shiver\textquotedblright{}.) Say it once more.
\end{tcolorbox}
\section[die Aufheiterung]{die Aufheiterung — the sunny spell, the bright interval}

\label{lesson:GE-C461-aufheiterung}



\begin{quote}
Before the new one: say the German for to drizzle, then the German for the shower.
\end{quote}

\subsection*{You'll want to know: die Aufheiterung}
\textbf{die Aufheiterung} — \textquotedblleft{}the sunny spell, the bright interval\textquotedblright{}.

\begin{quote}
Für den Abend meldet der Wetterbericht Aufheiterungen.
\end{quote}

\begin{quote}
\emph{The weather report forecasts sunny spells for the evening.}
\end{quote}

Plural: \emph{die Aufheiterungen}. A weather report word.

\begin{cousinweb}
\emph{Aufheiterung} is built on \emph{aufheitern} 'to brighten up', from \emph{heiter} 'bright, clear', an inherited Germanic adjective. EtymDB does not list the noun.
\end{cousinweb}

\begin{grammarlens}[title={What you've built}]
One more word for telling, explaining and describing at B1.
\end{grammarlens}

\subsection*{Your turn}
\begin{itemize}
\raggedright
  \item \emph{Say it:} \emph{die Aufheiterung}
  \item \emph{Say it:} \emph{die Aufheiterung}, once more
  \item \emph{Say it:} a sentence of your own with \emph{die Aufheiterung}
  \item \emph{Recall:} say the German for to drizzle, then the German for the shower, then say \emph{die Aufheiterung} again
\end{itemize}

\begin{tcolorbox}[breakable,colback=teal!4,colframe=teal!35!black,title={Before you move on}]
What does die Aufheiterung mean? (\textquotedblleft{}The sunny spell, the bright interval\textquotedblright{}.) Say it once more.
\end{tcolorbox}
\section[aufklaren]{aufklaren — to clear up (of the sky)}

\label{lesson:GE-C461-aufklaren}



\begin{quote}
Before the new one: say the German for the shower, then the German for the sunny spell.
\end{quote}

\subsection*{You'll want to know: aufklaren}
\textbf{aufklaren} — \textquotedblleft{}to clear up (of the sky)\textquotedblright{}.

\begin{quote}
Nach dem Gewitter klart der Himmel schnell auf.
\end{quote}

\begin{quote}
\emph{After the thunderstorm the sky quickly clears.}
\end{quote}

Separable: \emph{es klart auf}, \emph{es hat aufgeklart}. Not the same as \emph{aufklären} 'to explain'.

\begin{cousinweb}
\emph{aufklaren} is built from \emph{auf} and \emph{klar} 'clear', which reached Middle High German as \emph{klār} through Middle Dutch from Latin \emph{clārus} 'bright, clear' (EtymDB; Lewis \& Short). \emph{Aufklaren} is usually described as a sailors' word for the sky.
\end{cousinweb}

\begin{grammarlens}[title={What you've built}]
One more word for telling, explaining and describing at B1.
\end{grammarlens}

\subsection*{Your turn}
\begin{itemize}
\raggedright
  \item \emph{Say it:} \emph{aufklaren}
  \item \emph{Say it:} \emph{aufklaren}, once more
  \item \emph{Say it:} a sentence of your own with \emph{aufklaren}
  \item \emph{Recall:} say the German for the shower, then the German for the sunny spell, then say \emph{aufklaren} again
\end{itemize}

\begin{tcolorbox}[breakable,colback=teal!4,colframe=teal!35!black,title={Before you move on}]
What does aufklaren mean? (\textquotedblleft{}To clear up (of the sky)\textquotedblright{}.) Say it once more.
\end{tcolorbox}
\section[(dialogue)]{Words from chapters 454–461 — a review}

\label{lesson:GE-R461-review-454-461}



\begin{quote}
Say the German for the hike and the trail.
\end{quote}

\subsection*{Your turn}
Say these aloud:

\begin{itemize}
\raggedright
  \item the hike, the trail, to climb, to climb down, to rest
  \item the waterfall, the marsh, the small pond, to splash, to flow
  \item the clearing, the bushes, the fern, the moss, to rustle
  \item the dune, the low tide, the high tide, the bay, the shell
  \item the squirrel, the hedgehog, the wild boar, the stag, to chirp
  \item the farm, the pasture, the barn, to graze, to milk
  \item the weather report, the weather situation, cloudless, overcast, to change suddenly
  \item the drizzle, to drizzle, the shower, the sunny spell, to clear up
\end{itemize}

\begin{tcolorbox}[breakable,colback=teal!4,colframe=teal!35!black,title={Before you move on}]
The hike? (\textbf{die Wanderung}.) The trail? (\textbf{der Wanderweg}.) To climb? (\textbf{besteigen}.) To climb down? (\textbf{absteigen}.) To rest? (\textbf{rasten}.) The waterfall? (\textbf{der Wasserfall}.) The marsh? (\textbf{der Sumpf}.) The small pond? (\textbf{der Tümpel}.) To splash? (\textbf{plätschern}.) To flow? (\textbf{fließen}.) The clearing? (\textbf{die Lichtung}.) The bushes? (\textbf{das Gebüsch}.) The fern? (\textbf{der Farn}.) The moss? (\textbf{das Moos}.) To rustle? (\textbf{rascheln}.) The dune? (\textbf{die Düne}.) The low tide? (\textbf{die Ebbe}.) The high tide? (\textbf{die Flut}.) The bay? (\textbf{die Bucht}.) The shell? (\textbf{die Muschel}.) The squirrel? (\textbf{das Eichhörnchen}.) The hedgehog? (\textbf{der Igel}.) The wild boar? (\textbf{das Wildschwein}.) The stag? (\textbf{der Hirsch}.) To chirp? (\textbf{zwitschern}.) The farm? (\textbf{der Bauernhof}.) The pasture? (\textbf{die Weide}.) The barn? (\textbf{die Scheune}.) To graze? (\textbf{grasen}.) To milk? (\textbf{melken}.) The weather report? (\textbf{der Wetterbericht}.) The weather situation? (\textbf{die Wetterlage}.) Cloudless? (\textbf{wolkenlos}.) Overcast? (\textbf{bedeckt}.) To change suddenly? (\textbf{umschlagen}.) The drizzle? (\textbf{der Nieselregen}.) To drizzle? (\textbf{nieseln}.) The shower? (\textbf{der Schauer}.) The sunny spell? (\textbf{die Aufheiterung}.) To clear up? (\textbf{aufklaren}.)
\end{tcolorbox}
